\section{Engineering Challenge and Requirements}
\label{sec:requirements}

\subsection{Setting}

This work was motivated by an industrial ATM publish--subscribe programme whose 112-requirement baseline for pre-deployment verification spanned six areas: model generation from engineering artefacts, scenario generation, a field-records telemetry store, analytical preparation, a node--relationship core model, and design verification. Although the baseline never uses the term ``digital twin'', its closed loop from engineering models to synthetic scenarios and field telemetry defines one. Table~\ref{tab:requirements} condenses that baseline into seven domain-independent requirements (R1--R7).

\begin{table}[t]
\centering
\caption{Requirements for an architectural digital twin, condensed from a 112-requirement industrial ATM baseline. Status follows \S\ref{sec:maturity}: \emph{built} (implemented and tested), \emph{partly built}, or \emph{designed}.}
\label{tab:requirements}
\small
\begin{tabularx}{\textwidth}{@{}l Y l l@{}}
\toprule
ID & Requirement & Where & Status \\
\midrule
R1 & \textbf{Single source of truth.} Every twin artefact (analysis view, simulator input, counterpart) derives from one model instance via defined transformations. & \S\ref{sec:metamodel}, \S\ref{sec:m2s} & built \\
R2 & \textbf{Semantic consistency.} All consumers of dependency semantics apply identical rules, verified mechanically. & \S\ref{sec:consistency} & built \\
R3 & \textbf{Substrate separation.} Simulators read structural relations only, never derived analysis edges, avoiding tautological scoring. & \S\ref{sec:metamodel} & built \\
R4 & \textbf{Executable counterpart.} When the real system is unavailable, the model compiles into a runnable middleware deployment under the same faults. & \S\ref{sec:m2s} & built \\
R5 & \textbf{Calibration from telemetry.} Pre-deployment parameter guesses are re-estimated from traces, explicitly reporting non-identifiable parameters. & \S\ref{sec:s2m} & built (batch) \\
R6 & \textbf{Measured fidelity.} Twin-to-counterpart agreement is measured against declared tolerances and reported without gating builds. & \S\ref{sec:fidelity} & built; studies run \\
R7 & \textbf{Drift detection.} Model--system divergence triggers re-calibration or engineer-reviewed model repair. & \S\ref{sec:livetwin} & partly built$^{a}$ \\
\bottomrule
\multicolumn{4}{@{}p{\textwidth}@{}}{\footnotesize $^{a}$Batch parameter/behavioural drift detection, EMA re-calibration and residual feedback are evaluated on perturbed Docker counterparts (\S\ref{sec:evaluation}); topology drift is designed.}\\
\end{tabularx}
\end{table}

\subsection{Why these requirements are hard together}

Taken together, R1--R7 create four architectural tensions:
\begin{itemize}[leftmargin=*,nosep]
  \item \textbf{R1 vs.\ R3:} One model instance must serve both analysis and simulation while enforcing a strict view boundary (\S\ref{sec:metamodel}).
  \item \textbf{R4 vs.\ R6:} Generating the counterpart from the model guarantees partial agreement by construction; R6 is meaningful only with an explicit independence analysis (\S\ref{sec:independence}).
  \item \textbf{R5 vs.\ R6:} Calibration and evaluation traces must be strictly disjoint by protocol (\S\ref{sec:protocol}).
  \item \textbf{R7 depends on R5 and R6:} Detecting drift presupposes a calibrated baseline and measured noise; accordingly we built R5--R6 first, evaluated batch drift/EMA next (\S\ref{sec:evaluation}), and left topology drift as designed.
\end{itemize}
